% !TeX root = ../../Book2.tex
% 纲要标题：### 15.1 根的幂零性
% 纲要位置：计划纲要.md，第 1245 行。

\section{根的幂零性}
\label{b2:sec:1501}

有限维性使左理想的升降链都停止，但“链停止”还不等于“某个幂为零”。本节把两步分开：先用半单商描述根滤过的各层，再在最后一小节将有限维性与非交换 Nakayama 引理合起来，证明根确实幂零。始终令 \(A\) 为域 \(k\) 上的有限维结合含幺代数，记 \(J=J(A)\)。

% 纲要要点（供目录核对）：
%
% * 有限维代数的 Jacobson 根
% * 半单商与根滤过
% * 幂零根和一般根的区别
% * 对有限维代数 \(A\) 的降链 \(J(A)^n\) 先取稳定项，再用第14.2节 Nakayama 引理推出其为零；不把 Artin 性本身当作根幂零性的全部证明
%

\subsection{有限维代数的 Jacobson 根}
\label{b2:subsec:1501:01}

回顾\cref{b2:def:jacobson-radical}，\(J(A)\) 是全部极大左理想的交，并且是双边理想。它零化每个简单左模。具体地，简单模 \(S\) 由任意非零 \(s\) 生成，映射 \(A\to S\)、\(a\mapsto as\) 的核是极大左理想；故 \(Js=0\)，对任意 \(s\) 都成立。反过来，若 \(x\) 零化所有简单模，则对每个极大左理想 \(L\)，将它用于 \(A/L\) 中的 \(1+L\)，得到 \(x\in L\)，所以 \(x\in J\)。

\ProseParagraphBreak
每个简单 \(A\) 模都是 \(A\) 的商，因而有限维。每个有限生成 \(A\) 模也有限维：有限个生成元给出满射 \(A^n\to M\)。反过来，有限个 \(k\) 基向量当然也是 \(A\) 模生成元。以下讨论有限生成模时，可以使用这个维数条件，但不会把 \(A\) 模生成元数与 \(k\) 维数混为一谈。

\begin{definition}\label{b2:def:nilpotent-ideal}
设 \(I\) 为环 \(R\) 的双边理想。对正整数 \(N\)，\(I^N\) 是全部 \(N\) 个 \(I\) 中元素的乘积所生成的加法子群。若某个正整数 \(N\) 使 \(I^N=0\)，则称 \(I\) 为\term[mi-ling-lixiang]{幂零理想}[nilpotent ideal]，满足该式的最小正整数称为其\term[miling-zhishu]{幂零指数}[nilpotency index]。
\end{definition}

\begin{proposition}\label{b2:prop:nilpotent-ideal-in-radical}
任意含幺环的幂零双边理想都包含在其 Jacobson 根中。
\end{proposition}
\begin{proof}
若 \(I^N=0\)，则每个 \(x\in I\) 都满足 \(x^N=0\)，故 \(I\) 是 nil 双边理想。\cref{b2:cor:nil-ideal-radical}已证明这种理想包含于 Jacobson 根；它允许各元素使用不同的幂次，因而也适用于这里的共同幂次 \(N\)。
\end{proof}

例如 \(A=k[t]/(t^m)\)、\(m\ge1\)，其理想 \((\bar t)\) 幂零，且商为域 \(k\)。它因此恰是 \(J(A)\)：包含关系由上命题得到，而商的根为零，再用\cref{b2:prop:radical-quotient}即得反向包含。这里的根不是全部零因子的任意集合，而是由理想及简单模刻画的特定对象。

\subsection{半单商与根滤过}
\label{b2:subsec:1501:02}

\begin{proposition}\label{b2:prop:finite-algebra-semisimple-quotient}
设 \(k\) 为域、\(A\) 为有限维结合含幺 \(k\) 代数，\(J=J(A)\) 为其 Jacobson 根。商代数 \(A/J\) 是半单代数。每个被 \(J\) 零化的有限维左 \(A\) 模都是半单模。
\end{proposition}
\begin{proof}
由\cref{b2:prop:radical-quotient}，\(J(A/J)=0\)。商代数的左正则模有限维，故具有有限长度；其极大子模正是极大左理想，交为零。应用\cref{b2:prop:finite-length-radical-zero}，可知这个正则模半单，所以 \(A/J\) 是半单代数。

若 \(JM=0\)，定义 \((a+J)m=am\)。由于 \(J\) 在 \(M\) 上为零，此定义与代表元无关，使 \(M\) 成为 \(A/J\) 模。它是某个有限自由 \((A/J)^n\) 的商。有限直和与商保持半单性，见\cref{b2:prop:semisimple-subquotients}，故 \(M\) 半单；其 \(A/J\) 子模与 \(A\) 子模是同一批子空间。
\end{proof}

由\cref{b2:thm:artin-wedderburn}，可以写
\begin{equation}\label{b2:eq:finite-algebra-semisimple-quotient}
A/J\cong\prod_{i=1}^s M_{n_i}(D_i),
\end{equation}
其中各 \(D_i\) 是有限维 \(k\) 除代数。暂不假设 \(k\) 代数闭，所以不能把 \(D_i\) 一律写成 \(k\)。

\ProseParagraphBreak
设 \(k\) 为域、\(A\) 为有限维 \(k\) 代数，\(J=J(A)\) 为其 Jacobson 根。\(J\) 的各个幂给出\term[gen-lvguo]{根滤过}[radical filtration]
\[
A\supseteq J\supseteq J^2\supseteq\cdots.
\]
每个 \(J^r/J^{r+1}\) 都被 \(J\) 零化，因为 \(J J^r=J^{r+1}\)，所以由\cref{b2:prop:finite-algebra-semisimple-quotient}，它作为左模半单。类似地，对任何有限维左模 \(M\)，商 \(J^rM/J^{r+1}M\) 也是半单模。各层半单并不说明原模是这些层的直和；层与层之间的扩张可以不分裂。

\ProseParagraphBreak
在 \(k[t]/(t^m)\) 中，\(J^r\) 的基为 \(\bar t^r,\ldots,\bar t^{m-1}\)，\(0\le r<m\)。每个非零层都是一维的简单模 \(k\)，但当 \(m>1\) 时正则模不是这些层的直和：\(\bar t\) 在任何这样的简单模直和上均作用为零，而在正则模上并非如此。

\subsection{幂零根与一般根}
\label{b2:subsec:1501:03}

对有限维代数，本节将证明 \(J\) 是幂零理想。这个结论比“\(J\) 中每个元素各自幂零”更强，因为同一个整数要使任意若干根元素的乘积都为零。它也不意味着所有幂零元素都在根中。

\ProseParagraphBreak
例如 \(M_2(k)\) 是半单代数，故根为零，但矩阵单位 \(E_{12}\ne0\) 满足 \(E_{12}^2=0\)。幂零的单个元素生成的双边理想未必幂零；这里 \(E_{21}E_{12}=E_{22}\)，所以它生成的双边理想已经含有非零幂等元。\cref{b2:prop:nilpotent-ideal-in-radical}的假设是整个双边理想幂零。

\ProseParagraphBreak
若去掉有限维性，Jacobson 根可以完全不幂零。例如形式幂级数环 \(k[[t]]\) 中，常数项非零的级数可以逐次解出逆级数的系数，因而可逆；常数项为零的级数不可能可逆。故非单位元恰为理想 \((t)\)，它是唯一极大理想，等于 Jacobson 根。但是 \(t^n\ne0\) 对所有 \(n\) 都成立，根的各个幂始终非零。下一小节的有限维条件正用来排除这种不会停止的严格降链。

\subsection{根幂零性的 Nakayama 证明}
\label{b2:subsec:1501:04}

\begin{theorem}\label{b2:thm:finite-algebra-radical-nilpotent}
设 \(k\) 为域、\(A\) 为有限维结合含幺 \(k\) 代数，\(J=J(A)\) 为其 Jacobson 根，则 \(J\) 幂零。若 \(A\ne0\)，可以取正整数 \(N\le\dim_k A\) 使 \(J^N=0\)。因此 \(J\) 是 \(A\) 中最大的幂零双边理想。
\end{theorem}
\begin{proof}
降链 \(A\supseteq J\supseteq J^2\supseteq\cdots\) 是有限维向量空间的降链，故存在 \(n\ge0\) 使 \(J^n=J^{n+1}\)。令 \(M=J^n\)，将它看成左 \(A\) 模。它有限维，因而有限生成；等式正是
\[
JM=J J^n=J^{n+1}=J^n=M.
\]
对理想 \(J\subseteq J(A)\) 和有限生成左模 \(M\) 使用\cref{b2:thm:noncommutative-nakayama}，得到 \(M=0\)。所以某个 \(J^n\) 为零，证明幂零性。

当 \(A\ne0\) 时，\(J\ne A\)，因为单位不属于任何真极大左理想。若在到达零以前某两个相邻幂相等，上面的同一论证会迫使它们为零，矛盾。因此从 \(A\) 到第一个零幂，每一步维数都严格下降，步数不超过 \(\dim_k A\)，得到所述界。最后，任意幂零双边理想由\cref{b2:prop:nilpotent-ideal-in-radical}包含于 \(J\)，而 \(J\) 本身已证幂零。
\end{proof}

证明用到了两个不同的事实：有限维性使某一项稳定，Nakayama 引理使这个稳定项为零。仅凭降链停止，只能得到 \(J^n=J^{n+1}\)，不能省去后一步。这个幂零性也将保证稍后提升幂等元的过程在有限步后结束。
